\chapter{第三章：转行路径与岗位分析}

理解行业和技术版图之后，真正决定你能不能顺利转行的，不是``AI 很热''这个事实，而是：\textbf{你应该投什么岗位、补什么能力、用多长时间补、做什么作品集最有效。}


很多转行失败，并不是因为不够努力，而是因为目标岗位选错了。前端去冲模型训练岗，后端只刷 Prompt，QA 忽略评估体系，DevOps 完全不碰 LLM API------这些都很常见。


这一章的目标，就是把 AI Agent 相关岗位拆清楚，并给出从不同背景切入的可执行路线。


\bigskip\hrule\bigskip

\section{3.1 六类最常见岗位：先把坐标系立住}

\subsection{AI Agent Engineer}

这是本书最核心的目标岗位，职责是构建能够完成任务的 Agent 系统，通常涉及：


\begin{itemize}[leftmargin=2em]
\item 工具调用；
\item 工作流编排；
\item RAG；
\item memory；
\item 安全与权限；
\item 评估与上线。
\end{itemize}

一句话：\textbf{让模型真正干活的人。}


\subsection{ML Engineer}

更偏模型训练、部署与推理优化，常见工作有：


\begin{itemize}[leftmargin=2em]
\item 训练或微调模型；
\item serving 与推理优化；
\item GPU 资源管理；
\item 数据流水线。
\end{itemize}

一句话：\textbf{让模型更强、更稳的人。}


\subsection{LLM Application Engineer}

这类岗位介于传统应用工程师和 Agent 工程师之间，主要做：


\begin{itemize}[leftmargin=2em]
\item LLM API 接入；
\item Prompt、摘要、分类、问答；
\item RAG 驱动的产品功能。
\end{itemize}

它是很多工程师进入 Agent 方向的过渡台阶。


\subsection{AI Infrastructure Engineer}

非常适合 DevOps / SRE 切入。重点包括：


\begin{itemize}[leftmargin=2em]
\item 推理集群；
\item 模型部署平台；
\item 监控与成本优化；
\item GPU 调度；
\item LLMOps / MLOps。
\end{itemize}

\subsection{Data Scientist}

重点在分析、实验和指标，而不是搭完整 Agent 系统。常做：


\begin{itemize}[leftmargin=2em]
\item 数据分析；
\item A/B 测试；
\item 实验设计；
\item 业务建模。
\end{itemize}

\subsection{Prompt Engineer}

到了 2026 年，纯 Prompt Engineer 正在被并入其他岗位。Prompt 仍然重要，但它越来越像一项基础技能，而不是单独职业。


\bigskip\hrule\bigskip

\section{3.2 一张最重要的岗位对比表}

\begingroup
\scriptsize
\setlength{\tabcolsep}{3pt}
\begin{longtable}{|p{0.13\textwidth}|p{0.13\textwidth}|p{0.13\textwidth}|p{0.13\textwidth}|p{0.13\textwidth}|p{0.13\textwidth}|p{0.13\textwidth}|}
\hline
岗位 & 主要职责 & 必备技能 & 美国年薪区间 & 中国年薪区间 & 典型雇主 & 日常工作 \\ \hline
AI Agent Engineer & 做 Agent 系统、工具链、RAG、评估、上线 & Python、API、RAG、框架、系统设计、观测 & 15万-32万美元，高级可到45万+ & 32万-100万人民币，头部可到200万+ & AI 创业公司、SaaS、金融科技、企业数字化团队 & 设计流程、接工具、做评估、上线 \\ \hline
ML Engineer & 训练/部署模型、推理优化 & Python、PyTorch、Serving、GPU、数据流水线 & 16万-30万美元 & 40万-120万人民币 & 大模型公司、自动驾驶、研究团队 & 训练、部署、调优 \\ \hline
LLM Application Engineer & 构建 LLM 应用功能 & Python/TS、LLM API、Prompt、RAG & 14万-24万美元 & 30万-80万人民币 & SaaS、知识管理、内容平台 & 接 API、做实验、上线功能 \\ \hline
AI Infrastructure Engineer & 平台、推理、稳定性、成本 & Kubernetes、Docker、GPU、云平台、监控 & 16万-28万美元 & 35万-100万人民币 & 云厂商、大型企业平台团队 & 部署、容量规划、可靠性建设 \\ \hline
Data Scientist & 分析与实验 & SQL、Python、统计、实验设计 & 12万-22万美元 & 25万-70万人民币 & 互联网、零售、金融 & 分析数据、跑实验 \\ \hline
Prompt Engineer* & 提示词与对话质量优化 & Prompt、评估、业务理解 & 多已并入其他岗 & 多已并入其他岗 & 少量咨询/内容团队 & 优化提示词、做标注 \\ \hline
\end{longtable}
\endgroup

\* 提示：Prompt Engineer 作为纯独立岗位正在减少，更常见的是并入 AI Agent、LLM 应用或产品岗位。


从这张表你应该得到一个结论：\textbf{对多数软件工程师而言，AI Agent Engineer 和 LLM Application Engineer 是最现实的切入点；对 DevOps/SRE 而言，AI Infrastructure Engineer 是黄金路径。}


\bigskip\hrule\bigskip

\section{3.3 不同背景的最佳转行路径}

\subsection{前端 -> AI}

前端最大的优势不是``会做页面''，而是会做\textbf{Agent UI 和人机协作界面}。2026 年很多 Agent 产品都需要：


\begin{itemize}[leftmargin=2em]
\item React / TypeScript；
\item 流式输出（streaming UI）；
\item 复杂状态管理；
\item 任务轨迹可视化；
\item 审批流和工作台交互。
\end{itemize}

推荐补课顺序：


\begin{enumerate}[leftmargin=2em]
\item Python 基础；
\item LLM API；
\item RAG；
\item 一个 Agent 框架；
\item 做 Agent 控制台项目。
\end{enumerate}

\subsection{后端 -> AI}

这是最强通用路径。因为 Agent 的大多数生产难点，本质上都是后端问题：


\begin{itemize}[leftmargin=2em]
\item API 编排；
\item 权限控制；
\item 状态管理；
\item 队列与重试；
\item 日志与可观测性；
\item 成本优化。
\end{itemize}

后端要补的主要是：Python、LLM 基础、RAG、Agent 框架、Evaluation。


\subsection{全栈 -> AI}

全栈最适合做端到端 Agent 产品，也是最容易做出``有 ownership 的作品集''的人群。你可以：


\begin{itemize}[leftmargin=2em]
\item 用 React/Next.js 做前端；
\item 用 Python/FastAPI 做后端；
\item 接模型、向量数据库、日志平台；
\item 做一个完整可演示的业务场景。
\end{itemize}

这条路线也很适合独立开发者（indie hacker）。


\subsection{DevOps / SRE -> AI Infra}

这是很多人低估的机会。你已有的能力------Kubernetes、CI/CD、监控、容量规划、安全、故障处理------在 AI 系统里全部有用。


重点补：


\begin{itemize}[leftmargin=2em]
\item 推理服务；
\item GPU 基础；
\item 吞吐/延迟优化；
\item LLMOps / MLOps；
\item 模型网关与成本核算。
\end{itemize}

\subsection{QA -> AI Evaluation}

QA 是被低估的转行路径。传统测试方法论在 Agent 时代非常值钱：


\begin{itemize}[leftmargin=2em]
\item 用例设计；
\item 回归测试；
\item 边界情况；
\item 缺陷分类；
\item 自动化测试。
\end{itemize}

迁移到 AI 方向后，就变成：


\begin{itemize}[leftmargin=2em]
\item prompt regression；
\item RAG benchmark；
\item 红队测试（red teaming）；
\item hallucination 测试；
\item 多轮任务成功率评估。
\end{itemize}

\bigskip\hrule\bigskip

\section{3.4 你已有的技能，如何映射到 Agent 岗位}

\begingroup
\scriptsize
\setlength{\tabcolsep}{3pt}
\begin{longtable}{|p{0.31\textwidth}|p{0.31\textwidth}|p{0.31\textwidth}|}
\hline
你已有的能力 & 在 Agent 岗位中的价值 & 还要补什么 \\ \hline
Web/API 开发 & 工具接入、服务编排 & LLM API、tool calling、streaming \\ \hline
数据库设计 & memory、日志、知识库元数据 & 向量检索、embedding、RAG \\ \hline
分布式系统 & 队列、重试、扩缩容 & Agent runtime、workflow \\ \hline
前端交互 & Agent 控制台、审批流、可视化 & AI UX、对话状态设计 \\ \hline
测试自动化 & 回归与质量保障 & LLM evaluation、红队测试 \\ \hline
DevOps/CI/CD & 部署、监控、可靠性 & GPU、推理服务、LLMOps \\ \hline
\end{longtable}
\endgroup

你可以把学习目标分成三层：


\subsection{第一层：必须会}

\begin{itemize}[leftmargin=2em]
\item Python
\item LLM API
\item Prompt 基础
\item RAG 基础
\item 一个 Agent 框架
\item Docker / Git / 基础部署
\end{itemize}

\subsection{第二层：拉开差距}

\begin{itemize}[leftmargin=2em]
\item Evaluation
\item Tracing / observability
\item Memory 系统
\item Guardrails
\item 成本优化
\item 多 Agent 编排
\end{itemize}

\subsection{第三层：冲高薪}

\begin{itemize}[leftmargin=2em]
\item 私有化部署
\item 模型路由
\item 推理服务优化
\item 企业 IAM 集成
\item MCP / A2A
\item 大规模生产经验
\end{itemize}

\bigskip\hrule\bigskip

\section{3.5 一个现实的 3-6 个月转行计划}

假设你在职，每周能投入 12 到 18 小时，下面的节奏对大多数人比较现实。


\subsection{第 1 个月：补基础}

目标：


\begin{itemize}[leftmargin=2em]
\item 学 Python；
\item 理解 token、context window、embedding；
\item 接至少 2 家模型 API；
\item 做 2 到 3 个小 demo。
\end{itemize}

建议产出：


\begin{itemize}[leftmargin=2em]
\item 问答 demo；
\item 摘要或分类 demo；
\item 一个简单工具型 Agent。
\end{itemize}

\subsection{第 2 个月：掌握 RAG 和 Agent 结构}

目标：


\begin{itemize}[leftmargin=2em]
\item 学文档切分、embedding、召回、重排；
\item 选一个框架做完整流程；
\item 理解 tool use 和状态机。
\end{itemize}

建议产出：


\begin{itemize}[leftmargin=2em]
\item 一个带知识库的问答系统；
\item 一个能调 2 到 3 个工具的单 Agent。
\end{itemize}

\subsection{第 3 个月：做出作品集项目}

目标：


\begin{itemize}[leftmargin=2em]
\item 选真实业务场景；
\item 做前后端完整链路；
\item 接日志、评估、权限控制；
\item 能讲清楚架构取舍。
\end{itemize}

推荐方向：


\begin{itemize}[leftmargin=2em]
\item 客服 Agent；
\item 财报分析 Agent；
\item 编码 Agent；
\item 浏览器任务 Agent。
\end{itemize}

\subsection{第 4-6 个月：冲面试}

重点：


\begin{itemize}[leftmargin=2em]
\item 补系统设计；
\item 做回归测试；
\item 录项目演示；
\item 按岗位 JD 补短板；
\item 开始投递和模拟面试。
\end{itemize}

这里特别重要的一点是：\textbf{把项目做稳，比再开一个新项目更有价值。}


\bigskip\hrule\bigskip

\section{3.6 转行最常犯的错误}

\begin{enumerate}[leftmargin=2em]
\item \textbf{把``会用 ChatGPT''当成职业技能}：用户能力不等于工程能力。
\item \textbf{疯狂追框架，不做项目}：框架会变，项目交付能力才会留下。
\item \textbf{只学 Prompt，不学工程}：上线后决定质量的往往是权限、缓存、日志、评估。
\item \textbf{只会 demo，不会量化结果}：要能说清楚成功率、延迟、成本、召回率。
\item \textbf{作品过于玩具化}：比起``会讲笑话的机器人''，企业更爱客服、分析、编码、评估系统。
\item \textbf{忽略评估与安全}：prompt injection、防越权、人工审批、回归测试，都是高频面试题。
\item \textbf{不看岗位 JD 就学习}：AI Infra、LLM App、Agent Platform、Evaluation 关注点完全不同。
\item \textbf{低估原有背景价值}：你不是从零开始，而是在已有工程栈上加一层智能系统能力。
\end{enumerate}

\bigskip\hrule\bigskip

\section{3.7 什么样的作品集最打动招聘方}

如果只能做 3 个项目，我建议优先做下面三类。


\subsection{项目 1：生产感强的客服 Agent}

至少包含：


\begin{itemize}[leftmargin=2em]
\item RAG 知识库；
\item 订单/工单工具调用；
\item 升级人工；
\item 审计日志；
\item 延迟与成本统计。
\end{itemize}

\subsection{项目 2：编码 Agent 或文档分析 Agent}

重点体现：


\begin{itemize}[leftmargin=2em]
\item 多步任务；
\item 工具链；
\item 回归测试；
\item 失败处理。
\end{itemize}

\subsection{项目 3：评估平台或 Guardrails 项目}

这类项目很能拉开差距，例如：


\begin{itemize}[leftmargin=2em]
\item Agent trace viewer；
\item 回归评估脚本；
\item prompt injection 测试套件；
\item 敏感操作审批流。
\end{itemize}

下面给一个简单的 Python 评估示例：


\begin{codeblock}
test_cases = [
    {"question": "退款多久到账？", "expected": "3 个工作日"},
    {"question": "订单未支付能退款吗？", "expected": "未支付"},
]


def evaluate(agent_fn):
    passed = 0
    for case in test_cases:
        answer = agent_fn(case["question"])
        if case["expected"] in answer:
            passed += 1
    return {"score": passed / len(test_cases), "passed": passed}
\end{codeblock}

这个例子很简单，但代表了一个关键思维：\textbf{不要只问\texttt{\texttt{@@P0@@}}能不能稳定通过一组任务''。}


\bigskip\hrule\bigskip

\section{3.8 社区、证书和资源怎么选}

\subsection{社区}

优先加入三类社区：


\begin{itemize}[leftmargin=2em]
\item 开源社区：GitHub、Hugging Face、LangChain/Dify 等项目讨论区；
\item 中文实践社区：微信群、飞书群、线下 Meetup；
\item 招聘信息密集社区：LinkedIn、X、国内招聘平台专题圈子。
\end{itemize}

目的不是看热闹，而是：


\begin{itemize}[leftmargin=2em]
\item 看真实案例；
\item 跟进岗位需求；
\item 找项目合作和内推。
\end{itemize}

\subsection{证书}

证书不能替代项目，但对某些方向有辅助价值：


\begin{itemize}[leftmargin=2em]
\item 云厂商 AI/ML 认证：适合 AI Infra；
\item Kubernetes / 云原生认证：适合部署与平台岗位；
\item 安全认证：适合 Guardrails 和企业平台。
\end{itemize}

更实用的策略通常是：\textbf{先做强项目，再补证书。}


\subsection{学习资源优先级}

建议顺序：


\begin{enumerate}[leftmargin=2em]
\item 官方文档；
\item 生产级开源项目；
\item 有完整架构说明的技术博客；
\item 面试复盘与岗位 JD；
\item 论文和 benchmark 报告。
\end{enumerate}

对转行者而言，先打通工程闭环，收益通常高于一开始就扎进论文细节。


\bigskip\hrule\bigskip

\section{3.9 最后给不同背景读者的一句话建议}

\begin{itemize}[leftmargin=2em]
\item 前端：重点押注 Agent UI、工作台、流式交互。
\item 后端：重点押注工具编排、RAG、评估和平台能力。
\item 全栈：尽快做端到端项目，展示 ownership。
\item DevOps/SRE：补 GPU 和推理服务，走 AI Infra 路线。
\item QA：把测试方法论迁移到 Evaluation 和红队测试。
\end{itemize}

转行 AI Agent，不是推倒重来，而是把你已有的软件工程能力，升级成``智能系统工程能力''。


\section{本章要点}

\begin{itemize}[leftmargin=2em]
\item 2026 年最常见的 AI 相关岗位包括：\textbf{AI Agent Engineer、ML Engineer、LLM Application Engineer、AI Infrastructure Engineer、Data Scientist、Prompt Engineer（逐步并入其他岗位）}。
\item 对大多数软件工程师来说，\textbf{AI Agent Engineer 和 LLM Application Engineer} 是最现实的切入口；对 DevOps/SRE，\textbf{AI Infrastructure Engineer} 是高匹配路径；对 QA，\textbf{AI Evaluation} 是被低估的机会。
\item 岗位选择必须结合原背景做映射，而不是盲目追最热名词。
\item 一个现实的转行周期通常是 \textbf{3 到 6 个月}，关键是做出有业务价值、工程深度、评估意识的作品集。
\item 最常见的错误包括：只会用模型、不做项目、只学 Prompt、忽视评估与安全、作品过于玩具化。
\item 最能打动招聘方的项目，往往具备：\textbf{真实业务场景、多工具调用、RAG、日志、评估、审批或安全设计}。
\end{itemize}

\section{延伸阅读}

\begin{itemize}[leftmargin=2em]
\item OpenAI、Anthropic、Google、Microsoft 的开发者与 Agent 文档
\item LangChain / LangGraph、CrewAI、Dify、AutoGen 官方教程
\item LangSmith、Arize、Weights \& Biases 关于 LLM evaluation 的资料
\item vLLM、TGI、KServe、Ray Serve 等推理与部署工具文档
\item 2025-2026 年 AI/LLM/Agent 岗位 JD 与薪酬报告
\end{itemize}
